% 第1章 1.3.1：一致收敛的 epsilon 带。
% (a) f(x)=0.5+0.3 sin(2 pi x)，eps=0.15，f_n=f+0.1 sin(9 pi x) 整体落在带内（最大偏差 0.1<0.15）。
% (b) x^4 与 x^12 在 [0,1] 上相对零函数的 eps=0.15 带：二者在 x=0.15^(1/4)≈0.622、0.15^(1/12)≈0.854 处越出。
\begin{tikzpicture}[font=\figurefont\footnotesize]
\pgfplotsset{every axis/.style={width=50mm,height=38mm,axis lines=left,
  x axis line style={draw=FigureStroke,line width=.6pt},y axis line style={draw=FigureStroke,line width=.6pt},
  tick style={draw=FigureStroke},tick label style={font=\figurefont\scriptsize,text=FigureInk},
  label style={font=\figurefont\scriptsize},
  xlabel style={at={(axis description cs:1,0)},anchor=north west,yshift=1mm},
  title style={font=\figurefont\footnotesize},samples=200}}
\begin{axis}[name=A,xmin=0,xmax=1,ymin=0,ymax=1.05,xtick={0,0.5,1},ytick={0,0.5,1},xlabel={$x$},title={(a) 落在同一条带内}]
\addplot[fill=FigureBlueFill,draw=none,domain=0:1] {0.65+0.3*sin(deg(2*pi*x))} \closedcycle;
\addplot[fill=white,draw=none,domain=0:1] {0.35+0.3*sin(deg(2*pi*x))} \closedcycle;
\addplot[draw=FigureBlue,line width=.6pt,dashed,domain=0:1] {0.65+0.3*sin(deg(2*pi*x))};
\addplot[draw=FigureBlue,line width=.6pt,dashed,domain=0:1] {0.35+0.3*sin(deg(2*pi*x))};
\addplot[draw=FigureBlue,line width=1.1pt,domain=0:1] {0.5+0.3*sin(deg(2*pi*x))};
\addplot[draw=FigureYellow,line width=1.1pt,domain=0:1] {0.5+0.3*sin(deg(2*pi*x))+0.1*sin(deg(9*pi*x))};
\node[font=\figurefont\scriptsize,text=FigureBlue] at (axis cs:0.18,0.93) {$f$};
\node[font=\figurefont\scriptsize,text=FigureYellow] at (axis cs:0.80,0.43) {$f_n$};
\end{axis}
\begin{axis}[at={(A.east)},anchor=west,xshift=12mm,xmin=0,xmax=1,ymin=-0.2,ymax=1.05,xtick={0,0.5,1},ytick={0,0.5,1},xlabel={$x$},title={(b) $x^n$ 在端点附近越出}]
\addplot[fill=FigureBlueFill,draw=none] coordinates {(0,-0.15) (1,-0.15) (1,0.15) (0,0.15)} -- cycle;
\addplot[draw=FigureBlue,line width=.6pt,dashed] coordinates {(0,0.15) (1,0.15)};
\addplot[draw=FigureBlue,line width=.6pt,dashed] coordinates {(0,-0.15) (1,-0.15)};
\addplot[draw=FigureGreen,line width=1.1pt,domain=0:1] {x^4};
\addplot[draw=FigureYellow,line width=1.1pt,domain=0:1,samples=300] {x^12};
\node[font=\figurefont\scriptsize,text=FigureGreen] at (axis cs:0.12,0.88) {$x^4$};
\node[font=\figurefont\scriptsize,text=FigureYellow] at (axis cs:0.12,0.68) {$x^{12}$};
\node[font=\figurefont\scriptsize,text=FigureBlue,anchor=west] at (axis cs:0.02,0.24) {$|y|\le\varepsilon$};
\end{axis}
\end{tikzpicture}
